\documentclass[a4paper, 10pt]{article}
% Per-subject config for this repo. See vendor/template/course-config.tex.example
% for what each macro is used for.

\newcommand{\CourseCode}{2110515}
\newcommand{\CourseName}{Intro to Quantum Computing}
\newcommand{\StudentID}{6733172621}
\newcommand{\StudentName}{Patthadon Phengpinij}
\newcommand{\DefaultCollaborators}{Claude (for\,\LaTeX\,styling and grammar checking)}

\usepackage{../vendor/template/CEDT-Assignment-style}

\begin{document}
\subject
\asgntitle{4}{}{Week 4}{\StudentID\ \StudentName}{\DefaultCollaborators}

% ================================================================================ %
%                                    Problem 01                                    %
% ================================================================================ %
\begin{problem}
\textbf{Find Eigenvalues using Colab}
\par Find the eigenvectors and eigenvalues of \( \sigma_x \) using colab.
You may use this code.
But please ask yourself, which basis are we using when we express \( \sigma_x \) in the following form?
\begin{codingbox}
sigmax = np.array([[0, 1], [1, 0]])
print(sigmax)

eigenvalues, eigenvectors = np.linalg.eig(sigmax)
print(eigenvalues)
print(eigenvectors)
\end{codingbox}
\end{problem}

\begin{solution}
	The solution goes here.
\end{solution}
% ================================================================================ %

% ================================================================================ %
%                                    Problem 02                                    %
% ================================================================================ %
\begin{problem}
\textbf{Construct Projection Operator}
\par Now, use hand calculation, derive the projection operators to project to the eigenstates of \( \sigma_z \).
\end{problem}

\begin{solution}
	The solution goes here.
\end{solution}
% ================================================================================ %

% ================================================================================ %
%                                    Problem 03                                    %
% ================================================================================ %
\begin{problem}
\textbf{Use Projection Operator to find Measurement probability}
\par Use the projector equation to find the probability of measuring spin up and down when it is applied to the eigenvectors of \( \sigma_x \).
\par \underline{\textit{Hints:}} you should get 50\% for all of them.
\end{problem}

\begin{solution}
	The solution goes here.
\end{solution}
% ================================================================================ %

% ================================================================================ %
%                                    Problem 04                                    %
% ================================================================================ %
\begin{problem}
\textbf{Vector normalization by hand and using python 2}
\par Check \textbf{Problem 3.} using co-lab.
\par \underline{\textit{Hints:}}
\begin{itemize}
	\item Use \texttt{np.transpose} for transpose.
	\item Use \texttt{np.dot} for multiplication, including outer product.
	\item Use \texttt{np.linalg.eig} to find eigenvalues.
\end{itemize}
\end{problem}

\begin{solution}
	The solution goes here.
\end{solution}
% ================================================================================ %

% ================================================================================ %
%                                    Problem 05                                    %
% ================================================================================ %
\begin{problem}
\textbf{Orthonormality of Unitary Matrix Row Vectors}
\par Prove that the row vectors of a unitary matrix are orthonormal.
\end{problem}

\begin{solution}
	The solution goes here.
\end{solution}
% ================================================================================ %

% ================================================================================ %
%                                    Problem 06                                    %
% ================================================================================ %
\begin{problem}
\textbf{Finding the Determinants of Pauli Matrices by Hand}
\par Find the determinants of \( \sigma_x \), \( \sigma_y \), and \( \sigma_z \).
What do you notice?
\end{problem}

\begin{solution}
	The solution goes here.
\end{solution}
% ================================================================================ %

% ================================================================================ %
%                                    Problem 07                                    %
% ================================================================================ %
\begin{problem}
\textbf{Finding the Determinants of Pauli Matrices using Colab}
\par Confirm your answers in \textbf{Problem 6.} using Colab.
This is a sample code.
\begin{codingbox}
M = np.array([[1, 0], [0, -1]])
print(np.linalg.det(M))
\end{codingbox}
\end{problem}

\begin{solution}
	The solution goes here.
\end{solution}
% ================================================================================ %

% ================================================================================ %
%                                    Problem 08                                    %
% ================================================================================ %
\begin{problem}
\textbf{Proof of Unitarity}
\par For an orthonormal set of basis, \( \ket{v_i} \), we have \( \braket{v_i \vert v_j} = \delta_{i,j} \).
We transform it using a matrix \( A \), i.e. they become \( A\ket{v_i} \).
Prove that if after the transformation, they are still orthonormal, then \( A \) is unitary.
\par \underline{\textit{Hint:}} Express the basis vectors in column form and each of them has only 1 entry.
Express \( A \) in its matrix form. The fact that after the transformation,
they are still orthonormal implies that the columns of \( A \) are orthonormal and thus \( A \) is unitary.
\end{problem}

\begin{solution}
	The solution goes here.
\end{solution}
% ================================================================================ %

% ================================================================================ %
%                                     Appendix                                     %
% ================================================================================ %

\myappendix{code/assignment-04-code.pdf}{Assignment 04 | Code Solution}

% ================================================================================ %

\end{document}